\documentclass[11pt]{article}
\usepackage[letterpaper,margin=1in]{geometry}
\usepackage{amsmath,amssymb,amsthm,mathtools}
\usepackage{graphicx}
\usepackage{xcolor}
\usepackage{enumitem}
\usepackage[colorlinks=true,linkcolor=blue!55!black,citecolor=blue!55!black,urlcolor=blue!55!black]{hyperref}

\newtheorem{theorem}{Theorem}
\newtheorem{lemma}[theorem]{Lemma}
\newtheorem{corollary}[theorem]{Corollary}
\newtheorem{proposition}[theorem]{Proposition}
\theoremstyle{remark}
\newtheorem{remark}[theorem]{Remark}

\newcommand{\B}{\mathbb B}
\newcommand{\C}{\mathbb C}
\newcommand{\D}{\mathbb D}
\newcommand{\M}{M_4(\C)}
\newcommand{\Span}{\operatorname{span}}
\newcommand{\Alg}{\operatorname{Alg}}
\newcommand{\diag}{\operatorname{diag}}

\title{Three Variables Suffice for Commuting Row Contractions on $\C^4$}
\author{Automated mathematics run \texttt{fa\_banach\_001}, lane 12}
\date{August 9, 2026}

\begin{document}
\maketitle

\begin{center}
\fbox{\parbox{0.92\textwidth}{
\textbf{Status.} Substantial partial result, likely valid.  We prove that the
optimal row-contractive von Neumann constants satisfy
$C_{d,4}=C_{3,4}$ for every $d\geq3$.  Thus the source question in matrix
dimension four reduces completely to deciding whether
$C_{2,4}<C_{3,4}$.  That final comparison remains open.
}}
\end{center}

\section{Source question and result}

Hartz, Richter, and Shalit prove a dimension-free-in-$d$ von Neumann
inequality for commuting row contractions of fixed matrix size
\cite{HRS}.  For $d,n\geq1$, let $C_{d,n}$ be the least constant such that
\begin{equation}\label{eq:defC}
 \|p(T)\|\leq C_{d,n}\sup_{z\in\B_d}|p(z)|
\end{equation}
for every polynomial $p\in\C[z_1,\ldots,z_d]$ and every commuting tuple
$T=(T_1,\ldots,T_d)$ on $\C^n$ satisfying
\begin{equation}\label{eq:row}
 \sum_{j=1}^dT_jT_j^*\leq I.
\end{equation}
In Remark 4.8 on page 12 they state that they cannot determine whether
\begin{equation}\label{eq:source-question}
 C_{2,n}<C_{d,n}
\end{equation}
for any $n>2$.

\begin{figure}[ht]
\centering
\includegraphics[width=0.98\textwidth]{figures/open_problem_crop.png}
\caption{Remark 4.8 on page 12 of arXiv:2109.08550.}
\end{figure}

Piekarz subsequently settled the first possible matrix size, proving
$C_{d,3}=C_{2,3}$ for all $d\geq2$ \cite{Piekarz}.  The following theorem
handles all variable dimensions beyond two in the next matrix size.

\begin{theorem}[Dimension-four stabilization]\label{thm:main}
For every $d\geq3$,
\begin{equation}\label{eq:main}
 \boxed{\ C_{d,4}=C_{3,4}.\ }
\end{equation}
Consequently, $\sup_d C_{d,4}=C_{3,4}$, and a strict dependence on $d$ in
matrix dimension four can only be the single jump
$C_{2,4}<C_{3,4}$.
\end{theorem}

\section{Idea of the proof}

Move one joint eigenvalue of a strict tuple to the origin with a ball
automorphism.  If there is more than one joint eigenvalue, primary
decomposition and Schur's sharp dimension bound for commutative matrix
algebras show that the transformed coordinate matrices span at most three
dimensions.  A unitary rotation of the variables then kills all later
coordinates.

If there is only one joint eigenvalue, all transformed coordinates are
commuting nilpotent $4\times4$ matrices.  Their span again has dimension at
most four.  The only apparent obstruction is the extremal four-dimensional
case.  Maximality then forces this span to be an algebra.  An elementary
centralizer argument in the strictly upper triangular algebra proves that
every product in this algebra is zero.

That exceptional square-zero case is harmless analytically.  The functional
calculus retains only a bounded holomorphic function's value and gradient at
the origin.  Schwarz--Pick bounds the gradient, and a one-variable disk
automorphism matches the same first jet.  The classical von Neumann
inequality therefore gives constant one.  In every case, either three
variables suffice or the tuple satisfies a stronger constant-one estimate.

\section{The maximal nilpotent algebra}

We use Schur's theorem in the form used in Lemma 4.5 of \cite{HRS}: a
commutative subalgebra of $M_r(\C)$ has dimension at most
$\lfloor r^2/4\rfloor+1$.

\begin{lemma}[Extremal nilpotent lemma]\label{lem:nilpotent}
Let $L\subset M_4(\C)$ be a four-dimensional linear space of pairwise
commuting nilpotent matrices.  Then
\[
 L^2=\{0\}.
\]
\end{lemma}

\begin{proof}
Let $\mathcal A=\Alg(I,L)$ be the unital algebra generated by $L$.  Schur's
bound gives $\dim\mathcal A\leq5$, while $\C I+L$ already has dimension five.
Thus $\mathcal A=\C I+L$.  If $x,y\in L$, then $xy$ is nilpotent because
$x$ and $y$ commute.  Write $xy=\alpha I+z$ with $z\in L$.  If
$\alpha\neq0$, then $\alpha I+z$ is invertible, a contradiction.  Hence
$xy\in L$, so $L$ is an algebra.

A commuting family of complex nilpotents is simultaneously strictly upper
triangularizable.  In a common basis write
\begin{equation}\label{eq:x-matrix}
x=\begin{pmatrix}
0&a&b&c\\0&0&d&e\\0&0&0&f\\0&0&0&0
\end{pmatrix}\in L.
\end{equation}
The matrix unit $E_{14}$ annihilates every strictly upper triangular
$4\times4$ matrix on both sides.  If $E_{14}\notin L$, then
$\C I+L+\C E_{14}$ is a six-dimensional commutative algebra, contrary to
Schur's bound.  Therefore $E_{14}\in L$.

Suppose that some $x\in L$ has $x^2\neq0$.  For
\[
y=\begin{pmatrix}
0&A&B&C\\0&0&D&E\\0&0&0&F\\0&0&0&0
\end{pmatrix},
\]
the equation $xy=yx$ is equivalent to
\begin{equation}\label{eq:commutation}
 aD=Ad,\qquad dF=Df,\qquad aE+bF=Ae+Bf.
\end{equation}

First assume $d\neq0$.  The first two equations determine $A$ and $F$ from
$D$.  Substitution into the third gives one nonzero linear relation among
$B,D,E$.  Indeed, if this relation vanished identically then $a=f=0$, which
would imply $x^2=0$.  With the free central coordinate $C$, the centralizer
of $x$ inside the strictly upper triangular matrices therefore has dimension
at most three.  It cannot contain the four-dimensional space $L$.

It remains that $d=0$.  Now
$x^2=(ae+bf)E_{14}\neq0$, so $(a,f)\neq(0,0)$, and the first two equations
in \eqref{eq:commutation} force $D=0$.  On $\C^4$ define the nondegenerate
alternating form
\[
 \Omega((a,b,e,f),(A,B,E,F))=aE+bF-Ae-Bf.
\]
Project the centralizer coordinates to $(A,B,E,F)$.  Its kernel on $L$ is
exactly $\C E_{14}$, so the image of $L$ has dimension three.  Pairwise
commutativity says that this image is $\Omega$-isotropic.  But an isotropic
subspace of a four-dimensional symplectic space has dimension at most two,
a contradiction.

We have proved $x^2=0$ for every $x\in L$.  Since $L$ is linear and
commutative,
$0=(x+y)^2=x^2+2xy+y^2$ gives $xy=0$ for all $x,y\in L$.
\end{proof}

\section{The spectral and variable reduction}

\begin{lemma}[Four-dimensional dichotomy]\label{lem:dichotomy}
Let $T=(T_1,\ldots,T_d)$ be a strict commuting row contraction on $\C^4$.
There are a ball automorphism $\theta\in\operatorname{Aut}(\B_d)$ and a
unitary rotation of the variables such that $S=\theta(T)$ is a strict
commuting row contraction and either
\begin{enumerate}[label=\textup{(\roman*)}]
\item $S_j=0$ for every $j>3$; or
\item $S_iS_j=0$ for every $1\leq i,j\leq d$.
\end{enumerate}
\end{lemma}

\begin{proof}
Choose $\lambda\in\sigma(T)$ and let $\theta$ take $\lambda$ to zero.  Ball
automorphisms preserve strict commuting row contractions.  For completeness,
for the standard involution $\varphi_a$ this follows from the defect identity
\begin{equation}\label{eq:defect}
\begin{aligned}
 I-\varphi_a(T)\varphi_a(T)^*
  ={}&(1-\|a\|^2)
  \left(I-\sum_j\overline{a_j}T_j\right)^{-1}
  \left(I-\sum_jT_jT_j^*\right)\\
  &\qquad\cdot
  \left(I-\sum_j a_jT_j^*\right)^{-1}>0.
\end{aligned}
\end{equation}
Unitary rotations preserve the same property.  The inverse automorphism
recovers $T$ from $S$, so the generated algebras, and hence the primary block
sizes, are preserved.

Let the joint generalized eigenspace blocks of $S$ have dimensions
$r_1,\ldots,r_m$, with $\sum r_k=4$.  On block $k$ write
\[
 S_i|_k=\lambda_i^{(k)}I+N_i^{(k)},
\]
where the $N_i^{(k)}$ are commuting nilpotents.  Schur's bound applied to
their generated unital algebra gives
\begin{equation}\label{eq:block-bound}
 \dim\Span_i\{N_i^{(k)}\}\leq\left\lfloor\frac{r_k^2}{4}\right\rfloor.
\end{equation}
One spectral point is zero, so the scalar block patterns span at most
$m-1$ dimensions.  Consequently,
\begin{equation}\label{eq:rank-bound}
 \dim\Span\{S_1,\ldots,S_d\}
 \leq (m-1)+\sum_{k=1}^m\left\lfloor\frac{r_k^2}{4}\right\rfloor.
\end{equation}
If $m\geq2$, the four possible partitions give respectively
\[
3+1:3,\qquad 2+2:3,\qquad 2+1+1:3,\qquad 1+1+1+1:3.
\]
Thus the coordinate span has dimension at most three.

If $m=1$, every $S_i$ is nilpotent and Schur's bound gives
$\dim\Span\{S_i\}\leq4$.  Dimension at most three gives (i).  In dimension
four, Lemma~\ref{lem:nilpotent} gives (ii).  Finally, whenever the coordinate
span has dimension at most three, choose an orthonormal basis of the
coefficient space adapted to the kernel of
$c\mapsto\sum c_iS_i$.  The associated unitary variable rotation kills all
coordinates after the third.
\end{proof}

\section{The square-zero case}

\begin{lemma}[First-jet von Neumann inequality]\label{lem:squarezero}
Let $S=(S_1,\ldots,S_d)$ be a commuting row contraction with
$S_iS_j=0$ for all $i,j$.  If $q:\B_d\to\D$ is holomorphic, then
\[
 \|q(S)\|\leq1.
\]
\end{lemma}

\begin{proof}
Put $a=q(0)$, $b=(\partial_1q(0),\ldots,\partial_dq(0))$, and
$r=\|b\|_2$.  The infinitesimal Schwarz--Pick lemma on the ball \cite{Rudin}
gives
\begin{equation}\label{eq:SP}
 r\leq1-|a|^2.
\end{equation}
All terms of degree at least two vanish in the functional calculus, hence
\begin{equation}\label{eq:firstjet}
 q(S)=aI+\sum_jb_jS_j.
\end{equation}
If $r=0$ there is nothing to prove.  Otherwise let
\[
 A=\sum_j\frac{b_j}{r}S_j,
 \qquad s=\frac{r}{1-|a|^2}\leq1.
\]
Row contractivity gives $\|A\|\leq1$, and $A^2=0$.  The disk Schur function
\[
 h(z)=\frac{a+sz}{1+\overline a sz}
\]
satisfies $h(0)=a$ and $h'(0)=r$.  Therefore
$h(A)=aI+rA=q(S)$, and the classical von Neumann inequality (see, e.g.,
\cite[Section 1]{HRS}) gives
$\|q(S)\|\leq\|h\|_{\infty,\D}\leq1$.  If $|a|=1$, the maximum principle
makes $q$ constant, which is the omitted degenerate case.
\end{proof}

\section{Proof of the stabilization theorem}

\begin{proof}[Proof of Theorem~\ref{thm:main}]
Fix $d\geq3$, a commuting row contraction $T$ on $\C^4$, and a polynomial
$p$.  Applying the argument first to $\rho T$ for $0<\rho<1$ and then
letting $\rho\uparrow1$, it suffices to treat strict tuples.

Apply Lemma~\ref{lem:dichotomy}, including its unitary rotation, and put
$S=\theta(T)$ and $q=p\circ\theta^{-1}$.  The superposition property of the
holomorphic functional calculus gives
\begin{equation}\label{eq:superposition}
 p(T)=q(S),\qquad \|q\|_{\infty,\B_d}=\|p\|_{\infty,\B_d}.
\end{equation}
The rational function $q$ is holomorphic on a neighborhood of
$\overline{\B_d}$.

In case (i), let $\iota:\C^3\to\C^d$ be the coordinate inclusion.  Since
$S_j=0$ for $j>3$, functional calculus gives
$q(S)=(q\circ\iota)(S_1,S_2,S_3)$.  Uniform polynomial approximation of
$q\circ\iota$ on $\overline{\B_3}$ and the definition of $C_{3,4}$ yield
\begin{align*}
 \|p(T)\|=\|q(S)\|
 &\leq C_{3,4}\|q\circ\iota\|_{\infty,\B_3}\\
 &\leq C_{3,4}\|p\|_{\infty,\B_d}.
\end{align*}
In case (ii), normalize $q$ by its supremum norm and apply
Lemma~\ref{lem:squarezero}; this gives the stronger bound
$\|p(T)\|\leq\|p\|_{\infty,\B_d}$.  Hence $C_{d,4}\leq C_{3,4}$.

Conversely, pad every three-tuple with zero coordinates and regard every
three-variable polynomial as a $d$-variable polynomial.  Its supremum on
$\B_d$ equals its supremum on $\B_3$, so $C_{d,4}\geq C_{3,4}$.  This proves
\eqref{eq:main}.  Since $C_{1,4}=1\leq C_{2,4}\leq C_{3,4}$, the assertion
about the supremum follows as well.
\end{proof}

\section{Verification, novelty boundary, and review focus}

The exact checker in \texttt{code/verify\_dimension\_four\_nilpotent\_algebra.py}
enumerates all $11{,}011$ four-dimensional subspaces of the six-dimensional
strictly upper triangular algebra over $\mathbb F_3$.  It finds one
commutative subalgebra and verifies that it is square-zero, ending with
\texttt{ALL FINITE-FIELD NILPOTENT-ALGEBRA CHECKS PASSED}.  This is an exact
finite stress test, not a proof over $\C$.

The literature check was bounded.  On August 9, 2026, exact and close arXiv
searches used $C_{d,4}$, $c_{d,4}$, ``$4\times4$ row contractive von
Neumann,'' and close variants.  The source and Piekarz papers were inspected;
no statement of Theorem~\ref{thm:main} was found.  This is a bounded novelty
signal, not an exhaustive certification.

Human review should prioritize the maximality-to-algebra step, both
centralizer cases, the block dimension count, the defect identity, and the
first-jet matching.  Recommended status after review:
\texttt{substantial\_partial\_result\_likely\_valid}.

\section{Limitation}

The theorem completes the variable-stabilization program for matrix size
four but does not answer the source question there: it remains possible that
\[
 C_{2,4}<C_{3,4}.
\]
No numerical value or sharp bound for either constant is claimed.

\begin{thebibliography}{9}

\bibitem{HRS}
M. Hartz, S. Richter, and O. M. Shalit,
\emph{von Neumann's inequality for row contractive matrix tuples},
Math. Z. \textbf{301} (2022), 3877--3894;
\href{https://arxiv.org/abs/2109.08550}{arXiv:2109.08550}.

\bibitem{Piekarz}
D. Piekarz,
\emph{The von Neumann inequality for $3\times3$ matrices in the unit
Euclidean ball},
Israel J. Math. (2025);
\href{https://arxiv.org/abs/2310.12908}{arXiv:2310.12908}.

\bibitem{Rudin}
W. Rudin,
\emph{Function Theory in the Unit Ball of $\C^n$},
Springer-Verlag, New York, 1980.

\end{thebibliography}

\end{document}
